\section{The local contribution of a curve}\label{sec:partition}

How much does the normal operator of an embedded curve tell us about
the contribution of its multiple covers? Even if every cover is
isolated in the normal directions, counting the unramified covers with
their determinant signs gives only the first possible source genus in
each degree. Higher genera allow branching, nodes and contracted
components. We compute their combined contribution for a curve with
positive maximal flat symplectic normal holonomy. Already in degree
two, changing one of two identical spin structures can make the entire
disconnected contribution vanish. The reason appears when the target
is cut at a nonseparating node. Reversing one spin gluing tensors the
boundary character with the sign representation. It exchanges a Young
diagram with its transpose, so only self-conjugate diagrams contribute
when the two ends are glued back together.

Let \(C\) be a smooth projective complex curve of genus \(g\ge2\), and
let \(\rho:\pi_1(C)\to{\rm Sp}(4,\RR)\) be the holonomy of a flat
symplectic four-plane bundle. We use the convention that \(\rho\) is
\emph{positive maximal} if its compatible complex rank-two bundle
has degree \(2g-2\). In particular the total space has first Chern
number zero on its zero section.
The harmonic reduction gives a holomorphic bundle \(E\) and a
nondegenerate symmetric pairing \(\gamma:E\otimes E\to K_C\).
Choose a theta characteristic \(a\), that is, a holomorphic square
root of \(K_C\). Untwisting the pairing gives the complex orthogonal
bundle \(F=E\otimes a^{-1}\), the underlying bundle of the Cayley
correspondence of Bradlow, Garc\'ia-Prada and Gothen \cite{bgg}. Its first Stiefel--Whitney class \(w_1(F)\) measures
whether its transition functions can be chosen in \({\rm SO}(2,\CC)\).
In that case its two isotropic lines give \(F=M\oplus M^{-1}\);
write \(s=|\deg M|\). For \(w_1(F)\ne0\) monodromy interchanges the
isotropic directions, so an antidiagonal action scaling the two
lines separately need not exist.

A Hermitian metric turns the symmetric pairing into an antilinear map
\(R_\gamma:E\to\Lambda^{0,1}C\otimes E\) by taking the metric dual and
conjugating the form factor. The positive normal operator is
\(D=\bar\partial_E+R_\gamma\). In a split model \(E=a_1\oplus a_2\)
with \(a_i^2=K_C\), the pairing is the sum of the two square maps and a
product metric gives two independent real spin operators. For
\(w_1(F)\ne0\), every unmarked deformation type has a representative
of this form with \(a_1\ne a_2\). Conversely, the direct sum of the two uniformizing spin lifts realizes each such pair. Write \(\epsilon\) for the
determinant sign of \(D\) on \(C\), with orientation continued from the
complex Dolbeault operator. On such a representative
\(\epsilon=(-1)^{h^0(a_1)+h^0(a_2)}\).
Isolation, reduction to these representatives and invariance under
normal rescaling are proved in
\cite[Lemma~14.1, Propositions~14.2 and~15.4, and Appendices~D--E]{WueUnramified}.

Let \(\overline{\mathcal M}_k(C,m)\) be the proper stack of connected
unmarked stable maps of arithmetic genus \(k\) and degree \(m>0\) to
this fixed \(C\). Over a map \(f\), the antilinear term of \(D_f\)
uses \(df\), and sections on the branches of a node must agree.
The kernel is zero on every nonconstant component by the spin energy
identity. Matching then kills the constants on every contracted
subcurve. Its oriented real cokernel \(V_D\) consequently has rank
\(4(k-1-m(g-1))\), the real virtual dimension of the stack.
The ordinary connected local invariant and its disconnected series are
\begin{equation}\label{eq:fixed-target}
 L_{m,k}=\int_{[\overline{\mathcal M}_k(C,m)]^{\rm vir}}e(V_D),
 \qquad
 Z_D(Q,v)=\exp\!\left(\sum_{m\ge1,\,k}L_{m,k}Q^m v^{2k-2}\right).
\end{equation}
The sum starts at \(k=1+m(g-1)\) by Riemann--Hurwitz on the
normalization. Every connected component counted by \(Z_D\) has
positive degree. At the least genus each connected cover contributes its determinant sign divided by its deck order. Giving each cover the positive reciprocal of its deck order instead would recover the ordinary unramified Hurwitz number.

Bryan and Pandharipande computed the residue invariant of a direct sum
of two holomorphic line bundles with opposite torus weights
\cite[Corollaries~7.2--7.3]{bp08}. Their answer is a sum over all
partitions of the covering degree. Our first formula identifies the
real normal Euler integral with their residue invariant in the
orientable case. For a nonorientable \(F\), the sewing of the normal
operators changes that sum.

For a box \(u\) in a Young diagram \(\lambda\), write \(h(u)\) for
one plus the number of boxes strictly to its right or below it in its
row or column. Put \(a(\lambda)=\sum_{(i,j)\in\lambda}(j-i)\), the
total content. With \(t=e^{iv}\), let \(X_j=2-t^j-t^{-j}\) and
\(b_0=g-1\). Thus \(X_j=4\sin^2(jv/2)\).

\begin{theorem}\label{thm:spin}
For the positive maximal normal operator above, in every degree and
source genus the disconnected fixed-target series is
\[
 Z_D(Q,v)=
 \begin{cases}
 \displaystyle\sum_{\lambda}(-1)^{s|\lambda|}Q^{|\lambda|}
 t^{s a(\lambda)}\prod_{u\in\lambda}X_{h(u)}^{b_0},
       &w_1(F)=0,\\[5pt]
 \displaystyle\sum_{\lambda=\lambda^\top}(\epsilon Q)^{|\lambda|}
 \prod_{u\in\lambda}X_{h(u)}^{b_0},&w_1(F)\ne0.
 \end{cases}
\]
The empty diagram contributes one. The same series gives the local
invariants of a marked horizontal section of a flat symplectic
four-plane germ with this derivative holonomy and the split horizontal
and compatible vertical almost-complex structures.
\end{theorem}

The obstruction to the second comparison occurs at the seam of a
single handle. Pinch a nonseparating circle on \(C\), and choose the
two spin structures to agree on the normalization and to differ in the
identification of the two end fibers. For a relative map, the contact
orders make a partition \(\mu\) of \(m\). The mixed spin Euler integral
acquires the sign \(\epsilon^m(-1)^{m-\ell(\mu)}\) relative to the
square for a repeated spin. The second factor is the sign of a
permutation of cycle type \(\mu\). In the character basis, tensoring
by sign interchanges a partition and its transpose. Gluing the handle
takes a trace. A partition contributes to that trace exactly when
it is fixed by transposition. The orientation calculation has to
include even contacts and contracted components to give this sign
on the full relative virtual class.

Self-conjugate diagrams also occur in the local Real Gromov--Witten
theory of Georgieva and Ionel \cite[Theorem~1.1 and Corollary~9.14]{gi}.
Their domain and target carry antiholomorphic involutions and the
maps intertwine them. For their canonically oriented Real
Calabi--Yau series \(R_G(q,u)\) on a connected symmetric curve of
genus \(G\), the formal expression in the second line is precisely
\(R_{2g-1}((-1)^{g-1}\epsilon Q,iv)\). Indeed the crosscap sign in
their formula is raised to the even power \(2g-2\), and each
hyperbolic sine acquires a factor of \(i\). Our moduli in
\eqref{eq:fixed-target} are ordinary complex stable maps; the
orientation change comes from the \emph{normal spin matching} over
one node. In particular, taking the logarithm here selects connected
ordinary domains. In the Real theory the logarithm also counts pairs
of components exchanged by the involution \cite[(2.16)]{gi}.

For a concrete consequence, put \(x=X_1\). A balanced genus-\(j\)
factor is the disconnected antidiagonal local series of a curve of
genus \(j\) with both normal lines of degree \(j-1\); denote it
by \(B_j(Q,v)\). Its logarithm starts with \(Qx^{j-1}\).
The \emph{elementary multiplicities} \(n_{d,j}\) of a compactly
realized isolated germ are the integers in the factorization
\(Z_D=\prod_{d,j} B_j(Q^d,v)^{n_{d,j}}\), as justified by
\cite[Proposition~3.1]{WueUnramified}. The balanced local factor is the elementary curve of Ionel and Parker under the comparison of \cite[\S3]{WueProduct}. These multiplicities record a genus of the elementary
curve and differ from the integrals at a specified source genus.

\begin{corollary}\label{cor:cubic}
For two distinct spin structures of the same parity on a genus-two
target, the connected degree-three series is
\(28x^3/3-6x^4+x^5\). Its nonzero degree-three elementary
multiplicities in increasing genus are
\[
 (n_{3,4},n_{3,5},n_{3,6},n_{3,7})=(-64,64,-20,2).
\]
\end{corollary}
\begin{proof}
There is no self-conjugate partition of two. In degree three there
is just \((2,1)\), with hooks \(3,1,1\). Thus the logarithm gives
\(X_3x^2+x^3/3\), with \(X_3=9x-6x^2+x^3\). The only proper
divisor of three is one. Subtracting the degree-three term of the
balanced genus-two logarithm,
\(220x^3/3-70x^4+21x^5-2x^6\), gives the four stated integers.
\end{proof}

The degree-three answer is the same for any two distinct even
spin structures on a given genus-two curve. It gives the integral
of a product, summed over all stable maps. To recover the individual
classes one must locate them on the moduli space. We identify the
cokernel of a single spin operator with a real form of the residue
space at the ramification points. This gives a signed divisor formula
on the surface of covers with two simple branch points. On the
four-dimensional closure of the cubic covers with four simple branch
points the individual spin corrections have square \(6\) and mixed
pairing \(-1/3\) for any two of the ten distinct even spins
(Theorem~\ref{thm:gram}). Their common support is detected by the
difference of two squared Szeg\H{o} kernels. These intersections
explain what the partition sum forgets about an individual square
root of \(K_C\).
